\chapter{Lactose Tolerance Test}

\section*{What Is This Test?}
This test evaluates the body's ability to digest lactose, the primary sugar found in milk and dairy products. It is used to diagnose lactose intolerance, which is caused by a deficiency of the lactase enzyme \cite{misselwitz2013lactose}.

\section*{Alternative Names}
Hydrogen breath test for lactose, lactose breath test, oral lactose tolerance test

\section*{Principle}
Two methods are used:
\begin{itemize}
    \item \textbf{Hydrogen Breath Test (Standard):} Undigested lactose passes into the colon, where bacteria ferment it, producing hydrogen gas. The hydrogen is absorbed into the blood and exhaled. An increase in breath hydrogen indicates lactose malabsorption \cite{gasbarrini2009methodology}.
    \item \textbf{Lactose Blood Test (Less Common):} Measures blood glucose levels after a lactose load. In healthy individuals, lactase breaks lactose into glucose, raising blood sugar. A failure of blood glucose to rise indicates malabsorption.
\end{itemize}

\section*{History}
Lactose malabsorption was clinically described in the 1960s. The hydrogen breath test was developed in the 1970s and quickly became the preferred diagnostic standard due to its non-invasive nature and superior sensitivity compared to the blood glucose tolerance test.

\section*{Why This Test Is Ordered}
This test is ordered if you experience abdominal cramping, bloating, flatulence, nausea, or watery diarrhea after consuming dairy products.

\section*{Primary vs Secondary Test}
This is a primary diagnostic test for lactose malabsorption.

\section*{How This Test Is Performed}
- **Hydrogen Breath Test:** You fast for 12 hours. You provide a baseline breath sample, drink a lactose solution, and then blow into a breath analyzer at 15- to 30-minute intervals for 2 to 3 hours \cite{hammer2022european}.
- **Blood Test:** Serial blood samples are drawn at baseline, 30, 60, and 120 minutes after drinking the lactose solution.

\section*{What Preparation Is Needed}
Fast for 8 to 12 hours before the test. Do not smoke or exercise vigorously for 2 hours prior, and avoid high-fiber foods the day before. Avoid taking antibiotics or undergoing colonoscopy for 4 weeks before the test.

\section*{Contraindications and Precautions}
\begin{itemize}
  \item The test should not be performed in infants under 2 years of age due to the risk of severe dehydration from induced diarrhea. It is contraindicated in patients with suspected hereditary fructose intolerance.
\end{itemize}
\section*{Risks}
During the test, patients who are lactose intolerant may experience temporary gastrointestinal symptoms such as gas, bloating, cramps, or diarrhea.

\section*{What Happens After the Test}
If the test is positive, your doctor will advise dietary modifications, including reducing dairy intake, choosing lactose-free products, or taking lactase enzyme supplements (e.g., Lactaid) with meals.

\section*{Understanding Results}

\subsection*{Below Normal}
- **Breath Test:** A normal (negative) result. An increase of less than 20 ppm of hydrogen over baseline indicates normal lactose digestion.
- **Blood Test:** An abnormal result. A blood glucose rise of less than 20 mg/dL over baseline indicates lactase deficiency.

\subsection*{Above Normal}
- **Breath Test:** An abnormal (positive) result. An increase of 20 ppm or more of hydrogen over baseline indicates lactose malabsorption.
- **Blood Test:** A normal result. A rise of 20 mg/dL or more of blood glucose indicates normal digestion.

\section*{Normal Results}
- Breath hydrogen rise < 20 ppm above baseline.
- Blood glucose rise >= 20 mg/dL above baseline.

\section*{Follow-up Tests}
A trial of dietary elimination is often used to correlate test results with clinical symptoms.

\section*{References}
\bibliographystyle{unsrtnat}
\bibliography{references}

\clearpage
\section*{Regulatory Status and Authorization Date}
This chapter describes a test category, not a specific commercial product. FDA, CDSCO, EU IVDR, and MHRA approval or registration dates therefore cannot be assigned without the assay manufacturer, product name, instrument, and intended use. For laboratory-developed tests, an individual agency approval date may not exist. See the Preface for the interpretation of regulatory status.

\clearpage
